\documentclass[10pt,twoside,a5paper]{scrreprt}
\usepackage{hyperref}
\usepackage{graphicx}
\usepackage[giveninits=true,backend=biber,sortcites=true,%
uniquename=false,%
sorting=nyt,style=ext-authoryear-comp,dashed=false,%
maxnames=1000,doi=false,isbn=false,url=false,articlein=false,date=year]{biblatex}
\usepackage{amsfonts}
%\usepackage{mathscr}
\usepackage{mathrsfs}
\usepackage{caption}

\usepackage{scrlayer-scrpage}
\chead{\DMG{} Annual Report 2025}

\usepackage{fontspec}
\setmainfont{Latin Modern Math}
%\setmainfont{LMSans12}
%\setmainfont{LMRoman10}
%

%\usepackage[margin=2cm]{geometry}

\usepackage{wallpaper}


\usepackage{highlightAuthors}
\addbibresource{report.bib}
\addbibresource{unpub.bib}
\nocite{*}

\addapplicants{%
Bandara,%
Beliakov,%
Díaz Millán,%
James,%
Laird,%
Pineda-Villavicencio,%
Ranaweera,%
Ugon,%
}

\defbibcheck{year}{
  \iffieldint{year}{\ifnumequal{\thefield{year}}{2025}{}{\skipentry}}{\skipentry}}

\newcommand{\DMG}{$\mathbf{\mathbb{D}: \mathscr{M}\to G}$}


\title{Annual Report}
\author{} \date{2025}
\usepackage{pdfpages}

% This is necessary to eliminate some pandoc commands.
\newcommand{\pandocbounded}[1]{#1}
\begin{document}
  \begin{titlepage}
  \includepdf{Annual_Report.pdf}
  \end{titlepage}
  \tableofcontents


  \chapter*{Activities}
  \addcontentsline{toc}{chapter}{Activities}

  %Below is a list of activities by the group, in chronological order.

  \section*{Minicourse on Hodge Theory in Nonsmooth
  Settings}
  \addcontentsline{toc}{section}{Minicourse on Hodge
  Theory in Nonsmooth Settings}

  On 18, 20, 22nd August, \href{https://lashi.org}{Lashi Bandara} gave
  an online minicourse at the Geometric Analysis and Related Fields
  seminar, organised by
  \href{https://people.iut.ac.ir/en/sajjadlakzian}{Sajjad Lakzian} at
  Isfahan University of Technology. Recordings are available
  \href{https://www.youtube.com/watch?v=wllJ2wRyldw&list=PLuKJ_JmLD0R8a_PKKLO_ISLIII97cLBce&ab_channel=Lashi\%27sMathsChannel}{here}.

  \section*{2024 Yearbook}
  \addcontentsline{toc}{section}{25 August 2025: 2024 Yearbook}

  The 2024 Mathematics Yearbook came out in August. Congratulations to all the
  students who contributed to it. A fantastic recognition of their hard
  work. You can downlad it
  \href{https://blogs.deakin.edu.au/maths/yearbook/}{here}.

  \begin{figure}\centering\includegraphics[width=1.0\linewidth,height=\textheight,keepaspectratio]{../../content/post/25-09-02-yearbook/featured.jpg}\caption*{2024 Mathematics Yearbook Cover Page, based on student work}\end{figure}

  \section*{Trill White attends Heidelberg forum}
  \addcontentsline{toc}{section}{02 September 2025: Trill White
  attends Heidelberg forum}

  \href{../author/trill-white}{Trill White}  attended the \href{https://www.heidelberg-laureate-forum.org/}{12-th
  Heidelberg forum} from the 14th to the 19th of September.

  The HLF forum invites 200 selected students from mathematics and computer science across the world to meet with Fields, Turing, Abel and Nevanlinna award winners. Three students from Australian institutions attended, from Melbourne University, UNSW and Trill from Deakin, the only one coming from a non-Go8 university.

  while at this meeting, she handed the 2024 Deakin Mathematics Yearbook to
  Sarah MacLeod, Head of Young Researchers Relations.

  \begin{figure}\centering\includegraphics[width=1.0\linewidth,height=\textheight,keepaspectratio]{../../content/post/25-09-26-trillyearbook/featured.jpeg}\caption*{Trill White giving a copy of the 2024 Mathematics Yearbook ot Sarah MacLeod at HLF}\end{figure}

  \section*{Brianna Laird attends AMSI Winter
  School}
  \addcontentsline{toc}{section}{Brianna Laird
  attends AMSI Winter School}

  Brianna Laird attended the
  \href{https://ws.amsi.org.au/}{AMSI Winter School} from the first to the 12th of September.

  The AMSI Winter School is a nation-wide initiative by AMSI to organise week-long courses to mathematics students across the country. It is primarily targeted at research students. Brianna was one of the first Deakin student for many years to attend.

  \section*{Launch of VicGAP}
  \addcontentsline{toc}{section}{Launch of VicGAP}

  The inaugural seminar of the \href{https://vicgap.org}{Victorian GAP (Geometry, Analysis, and PDE)} network took place on 09/09/2025 at La Trobe University,
  City Campus.

  VicGAP is an initiative led by Lashi Bandara, who, together with Owen Dearicott (Latrobe) and Justin Forlano (Monash) to create a collaboration network for all Victorian mathematicians working in analysis, PDE and Geometry.

  More at: \href{https://vicgap.org/post/25-09-20/}{https://vicgap.org/post/25-09-20/}

   \section*{Guillermo gave a presentation at TU Berlin}
  \addcontentsline{toc}{section}{Guillermo gave a presentation at TU Berlin}

  \href{../authors/Guillermo-Pineda-Villavicencio}{Guillermo Pineda
  Villavicencio} gave a seminar at TU Berlin on the 17th of September.

  \textbf{Title:} Lower bound theorems on the face numbers of polytopes
  with few vertices

  \textbf{Abstract:} We study lower bound theorems for the number of
  \(k\)-faces (\(1\le k\le d-2\)) of a \(d\)-dimensional polytope \(P\)
  (or \emph{\(d\)-polytope}) with up to \(3d-1\) vertices. Earlier
  results include the following: Xue (2021) established the case of
  polytopes with at most \(2d\) vertices; Xue (2022) and
  Pineda-Villavicencio and Yost (2022) independently settled the case of
  \(2d+1\) vertices; and Pineda-Villavicencio, Tritama, and Yost (2024)
  extended this to \(2d+2\) vertices.

  We present a recent lower bound theorem covering \(d\)-polytopes with
  up to \(3d-1\) vertices. If \(P\) has exactly \(d+2\) facets and at
  least \(2d+\ell\) vertices, the lower bound is tight for certain
  combinations of \(d\) and \(\ell\). When \(P\) has at least \(d+3\)
  facets, the lower bound is always tight, with equality for some
  \(1\le k\le d-2\) only when \(P\) has precisely \(d+3\) facets. These
  results confirm a conjecture of Pineda-Villavicencio (2024).

  We outline the main ideas and methods underlying the proof and
  describe some minimisers for each number of vertices between \(d+2\)
  and \(3d-1\).

  We conclude with a discussion of what happens beyond \(3d-1\)
  vertices---where the picture becomes more intricate.

  This is joint work with Jie Wang.


  \section*{Iman Shekarriz visits Krakow, Leoben and
  Vienna}
  \addcontentsline{toc}{section}{Iman Shekarriz
  visits Krakow, Leoben and Vienna}

  \href{../author/iman-shekarriz}{Iman Shekarriz} visited the Technical
  University of Leoben from September 18th to 22nd. There, he met with
  \href{http://imrich.at/}{Wilfried Imrich} to discuss the theoretical
  advancement of soft happy colouring. Afterwards, they travelled
  together to at AGH university, Krakow, Poland to attend
  \href{https://www.10ccgt.agh.edu.pl/}{10CCGT}, a large gathering of
  graph theorists from all over the world. He presented a talk titled
  ``Soft Happy Colouring''. Upon returning to Austria, Iman stayed at
  the Vienna University of Technology to visit
  \href{https://svozil.github.io/publications/}{Karl Svozil} for some
  days. There, they discussed quantum chromatic contextuality.

  \begin{figure}\centering\includegraphics[width=1.0\linewidth,height=\textheight,keepaspectratio]{../../content/post/25-10-17-imantrip/featured.jpg}\caption*{Iman's collaboration mention by a speaker at the conference}\end{figure}

  \section*{Julien visits the University of Western
  Sydney}
  \addcontentsline{toc}{section}{Julien visits the
  University of Western Sydney}

  \href{../author/julien-ugon}{Julien Ugon} was kindly invited by Mehdi
  Tavakol at the University of Western Sydney to work on their joing
  project. While there, Julien gave a seminar on the teaching approach
  used in the school of IT at Deakin university, titled ``Rekindling
  students' engagement in mathematics through co-constructive
  practice''. The slides can be found
  \href{https://julien-ugon.gitlab.io/UWS-Presentation/Teaching/Presentation.html}{here}.

  \section*{AMSI Summer Research Scholarship
  success}
  \addcontentsline{toc}{section}{AMSI Summer
  Research Scholarship success}

  \href{../author/trill-white}{Trill White} and Hoa
  My Huy Lim received
  \href{https://srs.amsi.org.au/}{AMSI Summer Research Scholarships}.

  Trill worked under the supervision of \href{../author/kelvin-li}{Kelvin Li} on her project ``Optimisation in BKZ: An analysis of
  dynamic block size'' and Hoa worked under the supervision of Sergey Polyakovskiy on her project ``Just-in-Time Batch Scheduling
  Subject to Batch Size and Packing Constraints''.

  \section*{2025 ADR Prizes}
  \addcontentsline{toc}{section}{2025 ADR Prizes}

  Three mathematical students were awarded a 2025 ADR prize, to work on the following projects:

  \begin{itemize}
  
  \item
    \href{../author/brianna-laird}{Brianna Laird} under the supervision
    of \href{../author/reinier-diaz-millan}{Reinier Díaz Millán} and
    \href{../author/julien-ugon}{Julien Ugon} on ``Generalising the
    Weber Problem via Bregman Divergences in Infinite Dimensional
    Spaces''.
  \item
    \href{../author/kate-suraev}{Kate Suraev} under the supervision of
    \href{../author/julien-ugon}{Julien Ugon} and
    \href{../author/lashi-bandara}{Lashi Bandara} on ``Formal
    Verification of Mathematical Proofs in Lean''.
  \item
    Bao Minh Tran, under the supervision of
    \href{../author/gleb-beliakov}{Gleb Beliakov} and
    \href{../author/simon-james}{Simon James} on ``Out-of-distribution
    data fitting''.
  \end{itemize}

  Congratulations to our three students and we look forward to their
  projects.

  \section*{Reinier and Julien attend the Optimisation Days at UNSW}
  \addcontentsline{toc}{section}{Reinier and Julien
  attend the Optimisation Days at UNSW}

  \href{../author/julien-ugon}{Julien Ugon} and
  \href{../author/reinier-diaz-millan}{Reinier Diáz Millán} attended a
  two days workshop on optimisation, the
  \href{https://www.unsw.edu.au/science/our-schools/maths/our-research/optimisation/opt_days}{Optimisation
  Days} at the university of New South Wales. Reinier's trip was kindly
  supported by Vera Roshchina.

  Reinier gave a presentation titled ``Frank--Wolfe algorithm for DC
  optimisation problem''.

  \section*{Lashi visits IIT-Madras, Sri
  Jayawardenepura, Peradeniya and Kelaniya}
  \addcontentsline{toc}{section}{13 November 2025: Lashi visits
  IIT-Madras, Sri Jayawardenepura, Peradeniya and Kelaniya}

  \href{../author/lashi-bandara}{Lashi Bandara} visited
  \href{https://math.iitm.ac.in//innerfaculty.php?fname=Aprameyan\%20P}{Aprameyan
  Parthasarathy} at \href{https://math.iitm.ac.in/}{IIT-Madras} on
  invitation from 6th to 10th October.

  In the IIT-Madras Maths Colloquium, he gave a talk titled ``Functional
  calculus, garnished with a little geometry''.

  IIT-Madras maintains a beautiful tradition of keeping a collection of
  their colloquium talks hand-written by the speaker into a book.

  \begin{figure}\centering\includegraphics[width=1.0\linewidth,height=\textheight,keepaspectratio]{../../content/post/25-11-13-lashisouthasiatour/IITM-book.jpeg}\caption*{A page of the visitor's book at IITM}\end{figure}

  Together with B. Sivashankar, Aprameyan's PhD student, they discussed
  and began to collaborate on a project on understanding the stability
  of eigenvalues under rough metric convergence.

  During his leave visiting extended family in Sri Lanka, Lashi was
  invited by
  \href{https://science.sjp.ac.lk/mat/dr-p-g-t-harshani-2}{Thanuja
  Paragoda} from the \href{https://science.sjp.ac.lk/mat}{University of
  Sri Jayawardenapura} to give a talk to their undergraduate students.
  He gave a talk titled ``Crushing Orange Peels and Drawing Maps: the
  story of Gaussian curvature'', organised by the
  \href{https://science.sjp.ac.lk/mat/mathematics-society/}{Mathematics
  Society} of USJ.

  He was also invited by
  \href{https://sci.pdn.ac.lk/maths/staff/Ashwini-Amarasinghe}{Ashwini
  Amarasinghe} and
  \href{https://sites.google.com/view/gayanajayasinghe/home}{Gayana
  Jayasinghe} from the \href{https://pdn.ac.lk}{University of
  Peradeniya} and to give talks to their undergraduate students.
  Peradeniya has two departments of mathematics - one in the
  \href{https://sci.pdn.ac.lk/maths}{Faculty of Science} and the other
  in the \href{https://eng.pdn.ac.lk/maths/}{Faculty of Enginnering}. He
  gave the same talk as for USJ for the Science Faculty students at
  Peradeniya, who are focused more on pure mathematics, organised by the
  \href{https://mathema.soc.pdn.ac.lk}{Mathematical Society} of
  Peradeniya.

  \begin{figure}\centering\includegraphics[width=1.0\linewidth,height=\textheight,keepaspectratio]{../../content/post/25-11-13-lashisouthasiatour/Pera-poster.jpeg}\caption*{A poster advertising Lashi's talk at the university of Peradeniya}\end{figure}

  Lashi gave a chalkboard talk and he was very happy with the beautiful
  boards!

  \begin{figure}\centering\includegraphics[width=1.0\linewidth,height=\textheight,keepaspectratio]{../../content/post/25-11-13-lashisouthasiatour/Pera-talk.jpeg}\caption*{Lashi in the middle of a chalkboard presentation (Peradeniya)}\end{figure}

  For the Engineering faculty students, he gave a talk titled ``Chucking
  matrices into functions: why, when and how?'', organised by the
  \href{https://eng.pdn.ac.lk/maths/emsociety/}{Engineering Mathematics
  Society}.

  Lashi also gave an invited research colloquium talk at the
  \href{https://science.kln.ac.lk/depts/maths/}{University of Kelaniya}
  on invitation by
  \href{https://science.kln.ac.lk/depts/maths/index.php/component/sppagebuilder/?view=page&id=53}{Shamon
  Almeida}.

  \section*{Gleb Beliakov among Australia's
  top 250 researchers in 2026}
  \addcontentsline{toc}{section}{Gleb Beliakov
  among Australia's top 250 researchers in 2026}

  \href{../author/gleb-beliakov}{Gleb Beliakov} was named in the
  Australian's
  \href{https://www.theaustralian.com.au/special-reports/research-magazine/the-2026-research-magazine-showcases-the-benefits-of-research/news-story/cd04714e9c76e43bbcd32ad0288e7470}{2026
  Research Magazine} as a top researcher in Engineering and Computer
  Science, in the field of Fuzzy Systems. Congratulations to Gleb for
  this tremendous achievement.

  \section*{New book by Gleb Beliakov and Simon
  James}
  \addcontentsline{toc}{section}{26 November 2025: New book by Gleb
  Beliakov and Simon James}

  \href{../author/gleb-beliakov}{Gleb} and
  \href{../author/simon-james}{Simon}'s new book,
  \href{https://doi.org/10.1007/978-3-031-97070-2}{Choquet Capacities
  and Fuzzy Integrals}, co-authored with Jian-Zhang Wu, has just
  appeared.

  \section*{Simon and Julien attend Meanjin Delta}
  \addcontentsline{toc}{section}{Simon and Julien
  attend Meanjin Delta}

  \href{../author/simon-james}{Simon James} and
  \href{../author/julien-ugon}{Julien Ugon} attended the
  \href{https://fyimaths.weebly.com/meanjin-delta-2025.html}{Meanjin
  Delta conference} on the 25 and 26th of November. Julien presented a
  talk titled ``Co-constructive mastery-based learning in mathematics''
  and Simon's presentation was titled ``Five years of the mathematics
  yearbook: a side quest in co- constructive learning''. We acknowledge
  funding from the faculty of SEBE that enabled us to attend this
  conference.

  \section*{DMG Group attended the AustMS meeting}
  \addcontentsline{toc}{section}{13 December 2025: DMG Group attended
  the AustMS meeting}

  This week a delegation of six staff members, two PhD students and six
  undergraduate students from the DMG attended the
  \href{https://www.austms2025.org.au}{AustMS meeting} at Latrobe
  University. We gave the following talks:

  \begin{itemize}
  \item
    \href{../../author/lashi-bandara}{Lashi Bandara} -- \emph{The space
    of rough Riemannian metrics}; \emph{Harmonic analysis and boundary
    trace for first-order elliptic operators over vector bundles} and
    \emph{Cutting and gluing via matching boundary conditions}.
  \item
    \href{../../author/reinier-diaz-millan}{Reinier Díaz Millán} --
    \emph{Frank-Wolfe algorithm for star-convex functions}.
  \item
    \href{../../author/brianna-laird}{Brianna Laird} -- \emph{A Steiner
    Tree Based Genetic Algorithm for Multi-layered
    Location-Allocation-Routing Problems}.
  \item
    \href{../../author/guillermo-pineda-villavicencio}{Guillermo Pineda
    Villavicencio} - \emph{Colouring of graphs of polytopes}.
  \item
    \href{../../author/iman-shekarriz}{Iman Shekarriz} -
    \emph{Metaheuristic methods for Soft Happy Colouring}.
  \item
    \href{../../author/julien-ugon}{Julien Ugon} --
    \emph{Proximal-point-like algorithms for abstract convex
    minimisation problems}.
  \end{itemize}

  The ever ubiquitous \href{../../author/lashi-bandara}{Lashi} also
  participate in the debate ``AI will destroy mathematics'\,', which was
  both entertaining and illuminating, and
  \href{../../author/lashi-bandara}{Lashi} and
  \href{../../author/trill-white}{Trill White} each presented art work
  at the mathematical art exhibition.

  \begin{figure}\centering\includegraphics[width=1.0\linewidth,height=\textheight,keepaspectratio]{../../content/post/25-12-13-AustMS/featured.jpg}\caption*{A group of researchers and students attending AustMS 2025}\end{figure}

  \chapter*{Events}
  \addcontentsline{toc}{chapter}{Events}

  \section*{08 August 2025: DUMG Talk by A/Prof Nadia Sukhorukova:
  Mathematical Optimisation behind Deep Learning}
  \addcontentsline{toc}{section}{08 August 2025: DUMG Talk by A/Prof
  Nadia Sukhorukova: Mathematical Optimisation behind Deep Learning}

  Join the inagural Deakin Undergraduate Maths Group seminar by A/Prof
  Nadia Sukhorukova from Swinburne University.

  Title: ``Mathematical Optimisation behind Deep Learning''

  \section*{22 August 2025: Double DMG Seminar by Iman Shekkariz and
  Reinier Díaz Millán}
  \addcontentsline{toc}{section}{22 August 2025: Double DMG Seminar
  by Iman Shekkariz and Reinier Díaz Millán}

  \textbf{Title:} Soft Happy Colouring and Community

  A talk by Iman Shekkariz

  \textbf{Abstract:} For a coloured graph G and 0 ≤ ρ ≤ 1, a vertex v is
  ρhappy if at least ρ deg(v) of its neighbours share its colour. The
  soft happy colouring problem seeks a colouring σ that extends a given
  precolouring and maximises the number of ρ-happy vertices {[}3{]}.
  This NP-hard problem is closely linked to community detection in
  graphs. For example, for a graph in the stochastic block model (SBM)
  and for suitable ρ, with high probability, complete soft happy
  colourings can be achieved by the planted community structure {[}1{]}.
  Moreover, for 0 ≤ ρ1 \textless{} ρ2 ≤ 1, complete ρ2 -happy colourings
  achieve higher detection accuracy than complete ρ1 -happy colourings,
  and when ρ surpasses a critical threshold, it is unlikely to find a
  complete ρ-happy colouring with near-equal class sizes {[}2{]}.
  Finally, we survey existing algorithms and propose novel heuristic,
  local search, evolutionary, metaheuristic, and matheuristic approaches
  that enhance solution quality for soft happy colouring.

  References

  [1] Mohammad H. Shekarriz, Dhananjay Thiruvady, Asef
  Nazari, and Rhyd Lewis. Soft happy colourings and community structure
  of networks. Computers \& Operations Research, 174:106893, 2025.

  [2] Mohammad H. Shekarriz, Dhananjay Thiruvady, Asef Nazari, and
  Wilfried Imrich. Local search improvements for soft happy colouring.
  Preprint available on ArXiv: 2506.19284, 2025 1

  [3] Peng Zhang and
  Angsheng Li. Algorithmic aspects of homophyly of networks. Theoretical
  Computer Science, 593:117--131, 2015.

  \textbf{Title:} Faces of convex sets and its relation with the
  generalised interior

  A talk by Reinier Diaz Millan

  \textbf{Abstract:} This presentation explores the relationship between
  the facial structure of convex sets and various notions of generalised
  interiors, particularly in general vector spaces. We review the
  fundamental concepts of faces and minimal faces, with a central focus
  on the intrinsic core (icr), a generalisation of the relative interior
  that is characterised by its connection to minimal faces. A key result
  is the decomposition of any convex set into a disjoint union of the
  intrinsic cores of its faces. Building on this, we introduce the face
  relative interior (fri), a new concept of interiority for convex sets
  in topological vector spaces. We establish its fundamental properties
  and its place within the hierarchy of interiority notions. The talk
  also covers conditions that guarantee the face relative interior is
  nonempty. Ultimately, this work provides a deeper understanding of the
  geometric structure of convex sets in infinite-dimensional settings by
  linking their facial geometry to these generalised interior points.

  \section*{05 September 2025: DMG Talk by Prof Alexander Kruger}
  \addcontentsline{toc}{section}{05 September 2025: DUMG Talk by Prof
  Alexander Kruger}

    \textbf{Title:} About the Extremal Principle: From Convex Analysis to Nonsmooth Analysis (Geometric Considerations)

    \textbf{By:} Prof Alexander Kruger, Faculty of Mathematics and Statistics, Ton Duc Thang University, Ho Chi Minh City, Vietnam

    \textbf{Abstract:} Since the extremal principle was introduced in 1980, it has proved to be one of the key tools in nonsmooth optimization and variational analysis, serving as a substitution for the classical convex separation theorem when the convexity assumptions are not satisfied. Several extensions of the extremality property of collections of sets have been introduced as well as several extensions of the extremal principle.

    In this talk, I am going to present and discuss extremality, local extremality, stationarity and approximate stationarity properties of collections of sets and the corresponding (extended) extremal principle.

    References

    Kruger, A.Y., Mordukhovich, B.S.: Extremal points and the Euler equation in nonsmooth optimization problems. Dokl. Akad. Nauk BSSR 24(8), 684–687 (1980)

    Kruger, A.Y.: Weak stationarity: eliminating the gap between necessary and sufficient conditions. Optimization 53(2), 147–164 (2004)

    Bui, H.T., Kruger A.Y.: About extensions of the extremal principle. Vietnam J. Math. 46(2), 215–242 (2018)

    Bui, H.T., Kruger A.Y.: Extremality, stationarity and generalized separation of collections of sets. J. Optimization Theory Appl. 182(1), 211–264 (2019)

    Cuong, N.D., Kruger, A.Y., Thao, N.H.: Extremality of families of sets. Optimization 73(12), 3593–3607 (2024)

    Cuong, N.D., Kruger, A.Y.: Generalized separation of collections of sets. arXiv: 2412.05336

  \begin{figure}\centering\includegraphics[width=1.0\linewidth,height=\textheight,keepaspectratio]{../../content/event/25-09-05-kruger/featured.jpg}\caption*{Alex Kruger and attendees during Alex's DMG presentation}\end{figure}

  \section*{09 September 2025: DUMG Talk by Prof Alexander Kruger}
  \addcontentsline{toc}{section}{09 September 2025: DUMG Talk by Prof
  Alexander Kruger}

    \textbf{Title:} From Classical to Nonsmooth Analysis: Differentiating Nondifferentiable Functions

    \textbf{By:} Prof Alexander Kruger, Faculty of Mathematics and Statistics, Ton Duc Thang University, Ho Chi Minh City, Vietnam

    \textbf{Abstract:}
      I am going to trace the development of the idea of derivative from classical through convex to nonsmooth analysis. The emphasis is going to be on calculus (constant multiple and sum rules) and applications of differentiation to optimality conditions (Fermat theorem and multiplier rules).

      Conventional and strict derivatives (gradients) and several types of subdifferentials and normal cones will be discussed.

      No prior knowledge of convex or nonsmooth analysis is expected.

      References

      Rockafellar, R. T., Wets, R. J.-B. Variational Analysis, Springer-Verlag, Berlin, 1998

      Mordukhovich, B. S. Variational Analysis and Generalized Differentiation, Springer-Verlag, Berlin, 2006

      Penot, J.-P. Calculus Without Derivatives, Springer-Verlag, New York, 2013

      Kruger, A. Y. Properties of generalized differentials. Siberian Math. J. 26 (1985), 822–832

      Kruger, A. Y. On Fréchet subdifferentials. J. Math. Sci. 116, 3 (2003), 3325–3358 


  \section*{10 November 2025: DUMG Talk by Dr Reinier Díaz Millán}
  \addcontentsline{toc}{section}{10 November 2025: DUMG Talk by Dr
  Reinier Díaz Millán}

  Dr Reinier did his PhD in Optimisation from the Federal University of
  Goiás, Brazil and his talk is on ``The World's Most Famous Theorem'':
  Hidden within the fabric of mathematics lies a result so universal
  that it appears across cultures, sciences, and even art. Its countless
  proofs span from ancient scrolls to presidential sketches, yet its
  essence remains as captivating as ever. This talk will not only
  revisit the mystery behind it but also reveal why, after thousands of
  years, it continues to challenge, inspire, and surprise.

  \section*{18 November 2025: A day of \ldots{} Analysis \& Geometry}
  \addcontentsline{toc}{section}{18 November 2025: A day of \ldots{}
  Analysis \& Geometry}

  \href{https://tinyurl.com/AnalysisGeometryDeakin}{A day of \ldots{}
  Analysis \& Geometry} will take place at Deakin Waterfront on 8th
  December 2026. It has been designated as a
  \href{https://austms2025.org.au/satellites}{Satellite Conference} to
  the \href{https://www.austms2025.org.au}{Australian Mathematics
  Society 2026 Annual Meeting}.

  \begin{figure}\centering\includegraphics[width=1.0\linewidth,height=\textheight,keepaspectratio]{../../content/event/25-11-18-agd/featured.jpg}\end{figure}

  \chapter*{Group Members}
\section*{Researchers}
\begin{itemize}
\item Fatima Ansarizadeh
\item Lashi Bandara
\item Gleb Beliakov
\item Reinier Díaz Millán
\item Atabak Elmi
\item Peter Huf
\item Simon James
\item Kelvin Li
\item Vicky Mak-Hau
\item Guillermo Pineda Villavicencio
\item Julien Ugon
\end{itemize}
\section*{PhD Students}
\begin{itemize}
\item Mohammad Ghasemi
\item Brianna Laird
\item Iman Shekarriz
\end{itemize}
\section*{Honour Students}
\begin{itemize}
\item Trill White
\end{itemize}
\section*{Undergraduate Students}
\begin{itemize}
\item Brandon Smith
\item Kate Suraev
\item Tuan Thanh Lu
\end{itemize}


  \chapter*{Publications}
  \addcontentsline{toc}{chapter}{Publications}

  \printbibliography[check=year,type=article,heading=subbibliography,title={Refereed Articles}]
  \printbibliography[type=unpublished,heading=subbibliography,title={Preprints}]
\end{document}
